% main_metropolis.tex
\documentclass[aspectratio=169]{beamer}

% --- Langue & encodage ---
\usepackage[utf8]{inputenc}
\usepackage[T1]{fontenc}
\usepackage[english,french]{babel}
\usepackage{lmodern}
\usepackage{microtype}

% --- Maths & tableaux ---
\usepackage{amsmath,amssymb}
\usepackage{booktabs}
\usepackage{graphicx}
\usepackage{changepage}

% --- Thème Metropolis ---
\usetheme[numbering=fraction,progressbar=frametitle]{metropolis} % metro
\setbeamertemplate{frame numbering}[fraction] % 1/20


\definecolor{BootcampBlueDark}{HTML}{1A3C68}  % Bleu foncé pour l'en-tête
\definecolor{TextBlack}{HTML}{000000}         % Noir pour le texte
\definecolor{AccentGold}{HTML}{DAA520}       % Pour les vecteurs inconnus (x)
\definecolor{Crimson}{HTML}{990000}          % Pour les résultats (b)

\definecolor{BootcampBlueDark}{HTML}{1A3C68}
\definecolor{AccentGold}{HTML}{DAA520}
\definecolor{MatrixGreen}{HTML}{228B22}
\definecolor{MatrixRed}{HTML}{B22222}

\setbeamercolor{block title example}{bg=BootcampBlueDark!90!black,fg=white}
\setbeamercolor{block body example}{bg=BootcampBlueDark!5!white,fg=black}

% === Titre des frames (bande épaisse) ===
\setbeamercolor{frametitle}{bg=BootcampBlueDark,fg=white}
\setbeamerfont{frametitle}{series=\bfseries,size=\large}

% �� Augmente ici les dimensions ht=... et dp=...
\setbeamertemplate{frametitle}{
  \nointerlineskip
  \begin{beamercolorbox}[wd=\paperwidth,ht=3.2ex,dp=1.4ex,leftskip=1em,rightskip=1em]{frametitle}
    \usebeamerfont{frametitle}\insertframetitle
  \end{beamercolorbox}
}

\newenvironment{definitionblock}[1]{%
  \par\vspace{0.4em}%
  % ----- Titre : barre épaisse + marges -----
  \begin{beamercolorbox}[
      wd=\textwidth,
      sep=0.4em,      % <-- ÉPAISSEUR DE LA BARRE (padding vertical du titre)
      leftskip=0.3em, % <-- marge gauche du texte "Definition"
      rightskip=0.6em,
      rounded=false,
      shadow=false
    ]{block title example}%
    \raisebox{0.15em}{\usebeamerfont{block title example}#1}%
  \end{beamercolorbox}%
  \vspace{-0.4em}% rapproche un peu le corps du titre (ajuste si tu veux)
  % ----- Corps : fond clair + padding -----
  \begin{beamercolorbox}[
      wd=\textwidth,
      sep=1em,        % <-- padding interne du corps
      leftskip=0.6em, % <-- marge interne latérale du corps
      rightskip=0.6em,
      rounded=false,
      shadow=false
    ]{block body example}%
    \usebeamerfont{block body example}%
}{%
  \end{beamercolorbox}%
  \par\vspace{0.6em}%
}

% === Couleurs globales du texte et fond ===
\setbeamercolor{normal text}{fg=TextBlack,bg=white}
\setbeamercolor{title}{fg=TextBlack}
\setbeamercolor{structure}{fg=BootcampBlueDark}
\setbeamercolor{progress bar}{fg=BootcampBlueDark,bg=gray!30}


% --- Infos cours ---
\title{Lecture 1 --- Mathematical Tools: Matrices and Linear Systems}
\author{Nathan Sauldubois}
\institute{NYU Finance and Risk Engineering}
\date{January 5, 2026}
% \titlegraphic{\includegraphics[height=1.2cm]{nyu_logo.pdf}} % optionnel


\metroset{
  sectionpage=progressbar,
  subsectionpage=progressbar
}

% Automatically show the transition pages (plain = no header/footer)
\AtBeginSection{
  \frame[plain]{\sectionpage}
}
\AtBeginSubsection{
  \frame[plain]{\subsectionpage}
}

\begin{document}

\maketitle

\begin{frame}{Outline}
  \tableofcontents[hideallsubsections]
\end{frame}

\section{Main Definitions}

  \begin{frame}{Outline}
    \tableofcontents[currentsection, hideothersubsections]
  \end{frame}

\subsection{Motivations \& Definition of Matrices}
\begin{frame}{Linear Systems}

\begin{block}{Example: From Linear System to Matrix Form}
Consider the system of equations:
\[
\begin{cases}
\textcolor{BootcampBlueDark}{2} \, \textcolor{AccentGold}{x_1} 
+ \textcolor{BootcampBlueDark}{3} \, \textcolor{AccentGold}{x_2} 
= \textcolor{Crimson}{5} \\
\textcolor{BootcampBlueDark}{1} \, \textcolor{AccentGold}{x_1}
- \textcolor{BootcampBlueDark}{4} \, \textcolor{AccentGold}{x_2} 
= \textcolor{Crimson}{-2}
\end{cases}
\]

We can write it compactly as:
\[
\textcolor{BootcampBlueDark}{A} 
\textcolor{AccentGold}{x} 
= 
\textcolor{Crimson}{b}
\]
where the definition of $\textcolor{BootcampBlueDark}{A} 
\textcolor{AccentGold}{x}$ will be given later and
\[
\textcolor{BootcampBlueDark}{
A =
\begin{pmatrix}
2 & 3 \\
1 & -4
\end{pmatrix}},
\quad
\textcolor{AccentGold}{
x =
\begin{pmatrix}
x_1 \\ x_2
\end{pmatrix}},
\quad
\textcolor{Crimson}{
b =
\begin{pmatrix}
5 \\ -2
\end{pmatrix}}.
\]
\end{block}
\end{frame}

\begin{frame}{Definition of Matrices}
\begin{definitionblock}{Definition}
A \textbf{matrix} is a rectangular array of real numbers, of size $n \times m$:
\vspace{-3mm}
$$
A = (a_{ij})_{1 \le i \le n,\, 1 \le j \le m}, \quad a_{ij} \in \mathbb{R}.
$$

\vspace{-3mm}
\noindent It represents a linear transformation from $\mathbb{R}^m$ to $\mathbb{R}^n$. 
It is denoted by $\mathcal{M}_{n, m}(\mathbb{R})$, where $n$ is the number of rows and $m$ is the number of columns. Whenever $m =n$, we simply write $\mathcal{M}_{n}(\mathbb{R})$.
\end{definitionblock}

\vspace{-3mm}
\pause

\begin{block}{Examples}
\vspace{-5mm}
\begin{itemize}
    \item $A = \begin{pmatrix} 1 & 2 & 3 \\ 4 & 5 & 6 \end{pmatrix}$ is a $2 \times 3$ matrix and $B = \begin{pmatrix} 7 & 8 \\ 9 & 10 \\ 11 & 12 \end{pmatrix}$ is a $3 \times 2$ matrix.
    \item $C = (5)$ is a $1 \times 1$ matrix $\Rightarrow$ equivalent to a scalar in $\mathbb{R}$.
    \item A $n \times 1$ matrix is a \textbf{column vector}; a $1 \times n$ matrix is a \textbf{row vector}.
\end{itemize}
\end{block}
\end{frame}

\subsection{Algebraic Operations on Matrices}
\begin{frame}{Operations on Matrices: Multiplication by a Real Number}

\begin{definitionblock}{Definition}
Let $\textcolor{BootcampBlueDark}{A} = (a_{ij})_{1 \le i \le m,\, 1 \le j \le n}$ be a matrix with real coefficients, 
and let $\textcolor{AccentGold}{\lambda} \in \mathbb{R}$.

\vspace{-5mm}

\[
\textcolor{AccentGold}{\lambda} \textcolor{BootcampBlueDark}{A}
= (\textcolor{AccentGold}{\lambda} a_{ij})_{1 \le i \le m,\, 1 \le j \le n}.
\]

\vspace{-2mm}
\noindent That is, each entry of $A$ is multiplied by the same real number $\lambda$.
\end{definitionblock}

\pause

\vspace{-5mm}

\begin{block}{Examples}

\vspace{-5mm}


\begin{itemize}
  \item For $\displaystyle 
    A = \begin{pmatrix} 1 & 2 \\ -1 & 0 \end{pmatrix}$,
    $
    3A = 
    \begin{pmatrix} 
    3 & 6 \\ -3 & 0 
    \end{pmatrix}
    $, 
    $
    -2A = 
    \begin{pmatrix} 
    -2 & -4 \\ 2 & 0 
    \end{pmatrix}.
    $
  \item Multiplying by $1$ leaves the matrix unchanged:
    $
    1 \cdot A = A.
    $
\end{itemize}
\end{block}

\pause
\vspace{-3mm}

\begin{alertblock}{Properties}
\begin{itemize}
  \item Distributivity: $(\lambda + \mu)A = \lambda A + \mu A$.
  \item Associativity: $(\lambda \mu) A = \lambda (\mu A)$.
\end{itemize}
\end{alertblock}

\end{frame}

\begin{frame}{Operations on Matrices: Addition of Two Matrices}

\begin{definitionblock}{Definition (Addition and Subtraction)}
Let $A=(a_{ij})$ and $B=(b_{ij})$ be two matrices in $\mathbb{R}^{m\times n}$ \textit{(of the same size)}.  

\vspace{-4mm}
\[
A + B := (a_{ij}+b_{ij})_{1\le i\le m,\,1\le j\le n}, 
\qquad
A - B := A + (-1)B.
\]

\vspace{-2.5mm}


\noindent Addition and subtraction are performed \textit{entry by entry}.
\end{definitionblock}

\pause

\begin{block}{Examples}
\vspace{0.5mm}
\small
\[
\begin{pmatrix} 1 & 2 \\ -1 & 0 \end{pmatrix}
+
\begin{pmatrix} 3 & -5 \\ 4 & 7 \end{pmatrix}
=
\begin{pmatrix} 4 & -3 \\ 3 & 7 \end{pmatrix},
\qquad
\begin{pmatrix} 1 & 2 \\ -1 & 0 \end{pmatrix}
-
\begin{pmatrix} 3 & -5 \\ 4 & 7 \end{pmatrix}
=
\begin{pmatrix} -2 & 7 \\ -5 & -7 \end{pmatrix}.
\]
\[
\text{Zero matrix: }\;
\mathbf{0}_{m\times n}=
\begin{pmatrix}
0&\cdots&0\\
\vdots&\ddots&\vdots\\
0&\cdots&0
\end{pmatrix},
\qquad
A+\mathbf{0}_{m\times n}=A.
\]
\end{block}

\end{frame}

\begin{frame}{Operations on Matrices: Addition of Two Matrices}

\begin{definitionblock}{Definition (Addition and Subtraction)}
Let $A=(a_{ij})$ and $B=(b_{ij})$ be two matrices in $\mathbb{R}^{m\times n}$ \textit{(of the same size)}.  

\vspace{-4mm}
\[
A + B := (a_{ij}+b_{ij})_{1\le i\le m,\,1\le j\le n}, 
\qquad
A - B := A + (-1)B.
\]

\vspace{-2.5mm}

\noindent Addition and subtraction are performed \textit{entry by entry}.
\end{definitionblock}

\vspace{-4mm}

\begin{alertblock}{Properties (for matrices of the same size)}
\vspace{-1mm}
\begin{itemize}
  \item Commutativity: $A+B=B+A$
  \vspace{-1.5mm}
  \item Associativity: $(A+B)+C=A+(B+C)$
  \vspace{-1.5mm}
  \item Neutral element: $\mathbf{0}_{m\times n}$
  \vspace{-1.5mm}
  \item Additive inverse: $A+(-A)=\mathbf{0}_{m\times n}$
  \vspace{-1.5mm}
  \item Compatibility: $\lambda (A + B) = \lambda A + \lambda B$.
\end{itemize}
\end{alertblock}

\pause
\vspace{-4mm}

\begin{block}{Vector Space Structure and Dimension}
With the usual scalar multiplication, $\mathcal{M}_{n, m}(\mathbb{R})$ is a \textbf{vector space} over $\mathbb{R}$ of \textbf{dimension} $mn$.  
A canonical basis is $\{E_{ij}\}_{1\le i\le m,\,1\le j\le n}$ where $(E_{ij})_{kl}=\delta_{ik}\delta_{jl}$.
\end{block}

\end{frame}

\begin{frame}{Examples: Addition and Scalar Multiplication of Matrices}

\begin{block}{}
Let
$
A =
\begin{pmatrix}
1 & -2 & 3 \\
0 & 4 & -1
\end{pmatrix},
\qquad
B =
\begin{pmatrix}
2 & 0 & -1 \\
-3 & 1 & 5
\end{pmatrix}.
$
Both are matrices in $\mathcal{M}_{2,3}(\mathbb{R})$.
\end{block}

\pause

\begin{block}{Operations}
\[
A - 2B =
\begin{pmatrix}
1 & -2 & 3 \\
0 & 4 & -1
\end{pmatrix}
-
2\begin{pmatrix}
2 & 0 & -1 \\
-3 & 1 & 5
\end{pmatrix}
=
\begin{pmatrix}
-3 & -2 & 5 \\
6 & 2 & -11
\end{pmatrix}.
\]
\pause
\[
A + B =
\begin{pmatrix}
3 & -2 & 2 \\
-3 & 5 & 4
\end{pmatrix},
\qquad
B + A =
\begin{pmatrix}
3 & -2 & 2 \\
-3 & 5 & 4
\end{pmatrix}.
\]
\end{block}

\pause

\begin{block}{Summation Notation}
The sum of several matrices of the same size is written
\[
\sum_{p=1}^k A_p = A_1 + A_2 + \cdots + A_k.
\]
This operation is associative and commutative, as long as all $A_p$ belong to $\mathcal{M}_{m,n}(\mathbb{R})$.
\end{block}

\end{frame}

\begin{frame}{Operations on Matrices: Multiplication}

\begin{definitionblock}{Definition and Dimension Compatibility}
Let $A = (a_{i, j})_{\substack{1 \leq i \leq n \\ 1 \leq j \leq m }}\in \mathcal{M}_{n,m}(\mathbb{R})$ 
and 
$B = (b_{i, j})_{\substack{1 \leq i \leq m \\ 1 \leq j \leq p }}\in \mathcal{M}_{m,p}(\mathbb{R})$.  
Their product $AB \in \mathcal{M}_{n,p}(\mathbb{R}) $ is defined by
\[
(AB)_{ij}=\sum_{k=1}^{m} a_{ik}\,b_{kj},\qquad 1\le i\le n,\;1\le j\le p,
\]
and has size $m\times p$.  
\textbf{Compatibility is required:} the number of columns of $A$ must equal the number of rows of $B$ ($m$ here).
\end{definitionblock}

\pause

\begin{exampleblock}{Output Dimensions (quick checks)}
\[
(2\times 3)\cdot(3\times 4)\Rightarrow 2\times 4,\qquad
(3\times 3)\cdot(3\times 1)\Rightarrow 3\times 1,\qquad
(1\times n)\cdot(n\times 1)\Rightarrow 1\times 1\;(\text{scalar}).
\]
\end{exampleblock}

\end{frame}

\begin{frame}{Operations on Matrices: Multiplication}

\begin{definitionblock}{Definition and Dimension Compatibility}
Let $A = (a_{i, j})_{\substack{1 \leq i \leq n \\ 1 \leq j \leq m }}\in \mathcal{M}_{n,m}(\mathbb{R})$ 
and 
$B = (b_{i, j})_{\substack{1 \leq i \leq m \\ 1 \leq j \leq p }}\in \mathcal{M}_{m,p}(\mathbb{R})$.  
Their product $AB \in \mathcal{M}_{n,p}(\mathbb{R}) $ is defined by

\vspace{-4mm}

\[
(AB)_{ij}=\sum_{k=1}^{m} a_{ik}\,b_{kj},\qquad 1\le i\le n,\;1\le j\le p.
\]

\vspace{-6mm}
\end{definitionblock}

\pause
\vspace{-9mm}
\begin{block}{}
\textbf{Dot product.} For $x \in \mathcal{M}_{1,n}(\mathbb{R})$ and $y\in \mathcal{M}_{n,1}(\mathbb{R})$,
$
x y \;=\;\sum_{i=1}^n x_i y_i
\quad\text{(size }1\times 1\text{)}$.


\textbf{Outer product (column by row).} For $x \in \mathcal{M}_{n,1}(\mathbb{R})$ and $y\in \mathcal{M}_{1,m}(\mathbb{R})$,
$
x y \in \mathcal{M}_{n,m} (\mathbb{R}) ,\qquad (x y)_{ij}=x_i\,y_j
$.
\end{block}

\pause

\begin{alertblock}{Properties}
\begin{itemize}
  \item Associative: $(AB)C=A(BC)$ when dimensions are compatible.
  \item Distributive: $A(B+C)=AB+AC$ and $(A+B)C=AC+BC$.
  \item \textbf{Not commutative in general:} $AB\neq BA$ (and sometimes $BA$ is not even defined).
\end{itemize}
\end{alertblock}

\end{frame}

\begin{frame}{Example $1$ — Product of a $2\times3$ and a $3\times2$ matrix}

\begin{block}{}
Define
\[
\textcolor{BootcampBlueDark}{A} =
\begin{pmatrix}
\textcolor{MatrixRed}{1} & \textcolor{MatrixGreen}{2} & \textcolor{AccentGold}{-1} \\
0 & 3 & 4
\end{pmatrix},
\qquad
\textcolor{AccentGold}{B} =
\begin{pmatrix}
\textcolor{MatrixRed}{2} & 1 \\
\textcolor{MatrixGreen}{-1} & 0 \\
\textcolor{AccentGold}{3} & 2
\end{pmatrix}.
\]
\pause
\[
AB =
\begin{pmatrix}
\textcolor{MatrixRed}{(1)(2)}+\textcolor{MatrixGreen}{(2)(-1)}+\textcolor{AccentGold}{(-1)(3)} &
\textcolor{MatrixRed}{(1)(1)}+\textcolor{MatrixGreen}{(2)(0)}+\textcolor{AccentGold}{(-1)(2)}\\[4pt]
(0)(2)+(3)(-1)+(4)(3) & (0)(1)+(3)(0)+(4)(2)
\end{pmatrix}
=
\begin{pmatrix}
-3 & -1 \\[4pt] 9 & 8
\end{pmatrix}.
\]

\pause

\textbf{Note:} $A$ is $2\times3$, 
$B$ is $3\times2$, so $AB$ is $2\times2$; 
$BA$ is defined but $3\times3$.
\end{block}

\end{frame}

\begin{frame}{Example $2$ — Square matrices and powers}

\begin{exampleblock}{Example — Matrix multiplication is not commutative}

\definecolor{BootcampBlueDark}{HTML}{1A3C68}
\definecolor{AccentGold}{HTML}{DAA520}
\definecolor{MatrixRed}{HTML}{B22222}

\vspace{1mm}
Let $
\textcolor{BootcampBlueDark}{A} =
\begin{pmatrix}
1 & 2\\
0 & 1
\end{pmatrix} \,,\,
\textcolor{AccentGold}{B} =
\begin{pmatrix}
0 & 1\\
1 & 0
\end{pmatrix}.
$
Compute both products:
\pause
\[
\textcolor{MatrixRed}{AB} =
\begin{pmatrix}
1 & 2\\
0 & 1
\end{pmatrix}
\begin{pmatrix}
0 & 1\\
1 & 0
\end{pmatrix}
=
\begin{pmatrix}
2 & 1\\
1 & 0
\end{pmatrix},
\qquad
\textcolor{MatrixRed}{BA} =
\begin{pmatrix}
0 & 1\\
1 & 0
\end{pmatrix}
\begin{pmatrix}
1 & 2\\
0 & 1
\end{pmatrix}
=
\begin{pmatrix}
0 & 1\\
1 & 2
\end{pmatrix}.
\]

\pause
\vspace{-5mm}
\[
\boxed{AB \neq BA}
\]
\textbf{Observation:} although both products are $2\times2$, swapping the order changes the result.
\end{exampleblock}

\vspace{-5mm}

\begin{block}{}
Let $A \in \mathcal{M}_n (\mathbb{R})$. The powers of a square matrix are defined recursively:

\vspace{-6mm}

\[
A^2 = A\cdot A,\quad A^3 = A^2\cdot A,\quad
A^n = A^{n-1}\cdot A.
\]
\pause
\textbf{Remark:} if $A$ and $B$ commute ($AB=BA$), then $(AB)^n=A^nB^n$.
\end{block}
\end{frame}



\begin{frame}{Example $3$ — Product of a Matrix and a Vector}
    Let 
\[
A =
\begin{pmatrix}
2 & -1 & 3\\
0 & 4 & 1
\end{pmatrix}
\in \mathbb{R}^{2\times 3},
\qquad
x =
\begin{pmatrix}
1\\
-2\\
0
\end{pmatrix}
\in \mathbb{R}^{3}.
\]

\pause

The product $A x$ is defined by
\[
(Ax)_i = \sum_{j=1}^{3} a_{ij} x_j,
\quad \text{so} \quad
A x =
\begin{pmatrix}
2(1) + (-1)(-2) + 3(0)\\[3pt]
0(1) + 4(-2) + 1(0)
\end{pmatrix}
=
\begin{pmatrix}
4\\[3pt]
-8
\end{pmatrix}.
\]
\end{frame}

\begin{frame}{Range of a Matrix}

\begin{definitionblock}{Range (Column Space)}
Let $A \in \mathcal{M}_{m,n}(\mathbb{R})$.  
The \textbf{Range} of $A$, denoted $\mathrm{Rg}(A)$, is defined as
\vspace{-3mm}
\[
\mathrm{Rg}(A)=\{Ax \mid x\in \mathbb{R}^n\} \subset \mathbb{R}^m.
\]
\vspace{-1.6mm}
It is the set of all vectors that can be written as a linear combination of the columns of $A$.
\end{definitionblock}

\pause

\begin{definitionblock}{Geometric Meaning}
The range represents the \textbf{subspace of the codomain} where the matrix $A$ can actually send vectors.  
Its dimension is called the \textbf{rank} of $A$.
\end{definitionblock}

\end{frame}


\begin{frame}{Kernel of a Matrix}

\begin{definitionblock}{Kernel (Null Space)}
Let $A \in \mathcal{M}_{m,n}(\mathbb{R})$.  
The \textbf{kernel} of $A$, denoted $\ker(A)$, is defined by
\vspace{-3mm}
\[
\ker(A)=\{x\in\mathbb{R}^n \mid Ax = 0\}.
\]
\vspace{-1.6mm}
It is the set of all vectors that are mapped to the zero vector by $A$.
\end{definitionblock}

\pause

\begin{definitionblock}{Geometric Meaning}
The kernel is a \textbf{subspace of the domain} and measures how much information is \textbf{collapsed} by $A$.  
Its dimension is called the \textbf{nullity} of $A$.
\end{definitionblock}

\end{frame}



\begin{frame}{Inverse of a Matrix}

\begin{definitionblock}{Identity Matrix}
The \textbf{identity matrix} of size $n$, denoted $I_n$, is defined by
\vspace{-3mm}
\[
(I_n)_{ij} =
\begin{cases}
1, & \text{if } i=j,\\
0, & \text{if } i\neq j.
\end{cases}
\vspace{-1.6mm}
\]
It satisfies $AI_n = I_mA = A$ for any $A \in \mathcal{M}_{m, n} (\mathbb{R})$.
\end{definitionblock}

\pause

\begin{definitionblock}{Inverse Matrix}
A square matrix $A \in  \mathcal{M}_{n} ( \mathbb{R}) $ is said to be \textbf{invertible}
if there exists a matrix $A^{-1} \in \mathcal{M}_{n} (\mathbb{R} )$ such that
\[
A\,A^{-1} = A^{-1}A = I_n.
\]
In this case, $A^{-1}$ is called the \textbf{inverse} of $A$.
\end{definitionblock}


\end{frame}

\begin{frame}{Inverse of a Matrix}
\begin{definitionblock}{Inverse Matrix}
A square matrix $A \in  \mathcal{M}_{n} ( \mathbb{R}) $ is said to be \textbf{invertible}
if there exists a matrix $A^{-1} \in \mathcal{M}_{n} (\mathbb{R} )$ such that
\[
A\,A^{-1} = A^{-1}A = I_n.
\]
In this case, $A^{-1}$ is called the \textbf{inverse} of $A$.
\end{definitionblock}

\begin{block}{Characterization by the Homogeneous System}
A square matrix $A$ is invertible \textbf{if and only if} the equation
\[
A x = 0
\]
has only the trivial solution $x = 0$.
\end{block}
\end{frame}

\begin{frame}{Invertibility, Image, and Kernel}

\begin{definitionblock}{Characterization of Invertibility}
Let $A \in \mathcal{M}_{n}(\mathbb{R})$.  
The matrix $A$ is \textbf{invertible} if and only if the following equivalent conditions hold:
\begin{itemize}
\item $\mathrm{Rg}(A)=\mathbb{R}^n$
\item $\ker(A)=\{0\}$ (the kernel contains only the zero vector),
\item the rank of $A$ is $n$.
\end{itemize}
\end{definitionblock}

\pause

\begin{definitionblock}{Intuition}
Invertibility means that $A$ neither \textbf{loses dimensions} nor \textbf{collapses} directions:
\begin{itemize}
\item full image  $\Rightarrow$ every vector can be reached,
\item trivial kernel $\Rightarrow$ different vectors are not sent to the same one.
\end{itemize}
Thus $A$ defines a \textbf{bijective linear map}.
\end{definitionblock}

\end{frame}

\begin{frame}{Inverse Matrix}
\begin{definitionblock}{Inverse Matrix}
A square matrix $A \in  \mathcal{M}_{n} ( \mathbb{R}) $ is said to be \textbf{invertible}
if there exists a matrix $A^{-1} \in \mathcal{M}_{n} (\mathbb{R} )$ such that
\[
A\,A^{-1} = A^{-1}A = I_n.
\]
In this case, $A^{-1}$ is called the \textbf{inverse} of $A$.
\end{definitionblock}
    
\begin{alertblock}{Key Properties of the Inverse (for invertible square matrices)}
Let $A,B \in \mathcal{M}_{n} (\mathbb{R} )$ be invertible matrices, and let $\lambda \in \mathbb{R}\setminus\{0\}$.
\vspace{-3mm}

\begin{itemize}
  \item \textbf{Uniqueness:} the inverse of $A$ is unique.
\vspace{-1mm}
  \item \textbf{Inverse of the inverse:}  $(A^{-1})^{-1} = A.$ 
\vspace{-1mm}
  \item \textbf{Inverse of a product:} $
  (AB)^{-1} = B^{-1}A^{-1}.$
  (Notice the reversed order.)
\vspace{-1mm}
  \item \textbf{Inverse of a scalar multiple:} $
  (\lambda A)^{-1} = \frac{1}{\lambda}A^{-1}.$
\end{itemize}
\end{alertblock}
\end{frame}


\subsection{Transpositions}

\begin{frame}{Transposition of a Matrix : Definition}
\begin{definitionblock}{Definition}
\vspace{-2mm}
Let $A = (a_{ij}) \in  \mathcal{M}_{m,n} ( \mathbb{R}) $.  
The \textbf{transpose} of $A$, denoted $A^\top$, is the $n\times m$ matrix defined by
\vspace{-7mm}
\[
(A^\top)_{ij} = a_{ji}, \quad \text{i.e. } 
A^\top =
\begin{pmatrix}
\text{--- } \text{row}_1(A) \text{ ---}\\
\vdots\\
\text{--- } \text{row}_m(A) \text{ ---}
\end{pmatrix}^\top
=
\text{columns of } A \text{ become rows of } A^\top.
\]
\vspace{-9mm}
\end{definitionblock}
\pause 
\vspace{-4mm}
\begin{exampleblock}{Example — Computing a Transpose Explicitly}
Let $A =
\begin{pmatrix}
\textcolor{MatrixRed}{1} & \textcolor{MatrixGreen}{3} & \textcolor{AccentGold}{-2}\\
\textcolor{MatrixRed}{0} & \textcolor{MatrixGreen}{4} & \textcolor{AccentGold}{5}
\end{pmatrix}
\in \mathbb{R}^{2\times 3}.
$ Then the transpose $A^\top$ is obtained by switching rows and columns:
$
A^\top =
\begin{pmatrix}
\textcolor{MatrixRed}{1} & \textcolor{MatrixRed}{0}\\
\textcolor{MatrixGreen}{3} & \textcolor{MatrixGreen}{4}\\
\textcolor{AccentGold}{-2} & \textcolor{AccentGold}{5}
\end{pmatrix}
\in \mathbb{R}^{3\times 2}.
$
\textcolor{MatrixRed}{Row 1} of $A$ become \textcolor{MatrixRed}{Column 1} of $A^\top$,  
\textcolor{MatrixGreen}{Row 2} of $A$ become \textcolor{MatrixGreen}{Column 2} of $A^\top$,  
and colors show which entries correspond directly.
\end{exampleblock}

\end{frame}

\begin{frame}{Transposition of a Matrix : dot Product}
\begin{definitionblock}{Definition}
Let $A = (a_{ij}) \in \mathcal{M}_{m,n} ( \mathbb{R}) $.  
The \textbf{transpose} of $A$, denoted $A^\top$, is the $n\times m$ matrix defined by
\vspace{-3mm}
\[
(A^\top)_{ij} = a_{ji}, \quad \text{i.e. } 
A^\top =
\begin{pmatrix}
\text{--- } \text{row}_1(A) \text{ ---}\\
\vdots\\
\text{--- } \text{row}_m(A) \text{ ---}
\end{pmatrix}^\top
=
\text{columns of } A \text{ become rows of } A^\top.
\]
\vspace{-9mm}
\end{definitionblock}

\begin{exampleblock}{Connection with the Dot Product}
\vspace{0.5mm}
For any vectors $x,y \in \mathbb{R}^d$ (seen as $d$ lines and $1$ row), we have \textcolor{red}{
$ x \cdot y = 
x^\top y = \sum_{i=1}^{d} x_i y_i,
$}
which is the usual \textbf{dot product} in $\mathbb{R}^d$.
Hence $x^\top y = y^\top x$.
\noindent 
For $A \in \mathcal{M}_{n, m} (\mathbb{R})$, $x \in \mathbb{R}^m$ and $y \in \mathbb{R}^m$ we have \textcolor{red}{ $ Ay \cdot x = A^\top x \cdot y $.}
\end{exampleblock}

\end{frame}

\begin{frame}{Rg, Kernel, and Orthogonality}

\begin{definitionblock}{Orthogonal Complement}
For a subspace $E \subset \mathbb{R}^n$, the \textbf{orthogonal complement} is
\[
E^{\perp}=\{x\in \mathbb{R}^n \mid \langle x,y\rangle = 0 \text{ for all } y\in E\}.
\]
\end{definitionblock}

\pause

\begin{definitionblock}{Fundamental Relations}
Let $A\in \mathcal{M}_{m,n}(\mathbb{R})$. Then
\[
\mathrm{Rg}(A^{T}) = (\ker A)^{\perp}
\qquad\text{and}\qquad
\mathrm{Rg}(A)= (\ker A^{T})^{\perp}.
\]
That is, the range of a matrix is the orthogonal complement of the kernel of its transpose.
\end{definitionblock}

\pause

\begin{definitionblock}{Meaning}
A vector $x$ is orthogonal to every vector in $\ker(A)$  
if and only if it lies in $\mathrm{Im}(A^{T})$.  

These identities express a deep link between
\textbf{row space, column space, and solutions of linear systems}.
\end{definitionblock}

\end{frame}

\begin{frame}{Transposition of a Matrix : Symmetric and Anti-symmetric Matrices }

\begin{definitionblock}{Symmetric and Skew-Symmetric Matrices}
A matrix $A\in \mathcal{M}_n ( \mathbb{R}) $ is:
\begin{itemize}
  \item \textbf{Symmetric} if $A^\top = A$.
  \item \textbf{Skew-symmetric (anti-symmetric)} if $A^\top = -A$.
\end{itemize}
\textbf{Examples:}
\[
\begin{pmatrix}
2 & 3\\[2pt]
3 & 5
\end{pmatrix}
\text{ is symmetric,}\quad
\begin{pmatrix}
0 & 2\\[2pt]
-2 & 0
\end{pmatrix}
\text{ is skew-symmetric.}
\]
\end{definitionblock}
\end{frame}

\begin{frame}{Positive (Semi)Definite Matrices and Order on Symmetric Matrices}

\begin{definitionblock}{Positive and Positive Semidefinite Matrices}
A symmetric matrix $A\in \mathcal{M}_n(\mathbb{R})$ is called:
\begin{itemize}
  \item \textbf{positive semidefinite} if 
  \[
  x^\top A x \ge 0 \quad \text{for all } x\in\mathbb{R}^n,
  \]
  \item \textbf{positive definite} if 
  \[
  x^\top A x > 0 \quad \text{for all } x\ne 0.
  \]
\end{itemize}
We then write $A\succeq 0$ (semidefinite) and $A\succ 0$ (definite).
\end{definitionblock}

\pause

\begin{definitionblock}{Order on Symmetric Matrices}
For symmetric matrices $A$ and $B$ of the same size, we write
\[
A \preceq B
\quad \Longleftrightarrow \quad
B-A \text{ is positive semidefinite}.
\]
This defines a \textbf{partial order} on the space of symmetric matrices.
\end{definitionblock}

\end{frame}

\begin{frame}{Properties of the Transposition}

\begin{alertblock}{Main Properties}
For all matrices $A,B$ (of compatible dimensions) and scalars $\lambda \in \mathbb{R}$:

\vspace{-2 mm}
\begin{itemize}
  \item $(A^\top)^\top = A$
\vspace{-1 mm}
  \item $(A + B)^\top = A^\top + B^\top$
\vspace{-1 mm}
  \item $(\lambda A)^\top = \lambda A^\top$
\vspace{-1 mm}
  \item $(AB)^\top = B^\top A^\top$
\vspace{-1 mm}
  \item If $A$ is invertible $
  (A^\top)^{-1} = (A^{-1})^\top. $
\end{itemize}
\end{alertblock}


\end{frame}

\begin{frame}{Orthogonal Group}

\begin{definitionblock}{Definition}
The \textbf{orthogonal group} of order $n$, denoted 
\vspace{-3mm}
\[
O(n) = \{\, Q \in \mathcal{M}_n (\mathbb{R}) \mid Q^\top Q = QQ^\top = I_n \,\},
\vspace{-3mm}
\]
is the set of all \textbf{orthogonal matrices}.
\end{definitionblock}

\pause
\vspace{-4mm}

\begin{block}{Equivalent Characterizations}
For $Q \in \mathcal{M}_n (\mathbb{R})$, the following are equivalent:
\begin{itemize}
\vspace{-3mm}
  \item $Q^\top Q = I_n$  \hfill (definition)
  \item $Q^{-1} = Q^\top$  \hfill (inverse equals transpose)
 \end{itemize}
\end{block}

\pause
\vspace{-3mm}

\begin{exampleblock}{Example}
\[
Q =
\begin{pmatrix}
0 & -1\\
1 & \phantom{-}0
\end{pmatrix}
\in O(2), \quad
Q^\top Q = 
\begin{pmatrix}
1 & 0\\
0 & 1
\end{pmatrix} = I_2.
\]
This matrix represents a $90^\circ$ rotation in the plane.
\end{exampleblock}

\end{frame}

\begin{frame}{Properties and Interpretation of $O(n)$}

\begin{alertblock}{Key Properties}
For $Q_1,Q_2\in O(n)$ and any $x,y\in\mathbb{R}^n$:
\vspace{-3mm}
\begin{itemize}
  \item $Q_1Q_2\in O(n)$ \hfill (closed under multiplication)
  \vspace{-2mm}
  \item $Q^{-1}=Q^\top\in O(n)$ \hfill (inverse and transpose remain orthogonal)
  \vspace{-2mm}
  \item $\det(Q)=\pm1$
  \vspace{-3mm}
\end{itemize}
Hence $O(n)$ is a \textbf{group} under matrix multiplication.
\end{alertblock}

\pause
\vspace{-3mm}

\begin{exampleblock}{Norm and Inner Product Preservation}
For any $Q\in O(n)$ and $x,y\in\mathbb{R}^n$:
\vspace{-3mm}
\[
\|Qx\|_2^2 = (Qx)^\top (Qx) = x^\top Q^\top Q x = x^\top x = \|x\|_2^2, \,\, \text{and} \,\,
\langle Qx, Qy\rangle = x^\top y.
\vspace{-4mm}
\]
\textbf{Interpretation:} Orthogonal matrices represent \textit{rotations or reflections}
that preserve lengths and angles.
\end{exampleblock}



\end{frame}

\subsection{Trace}

\begin{frame}{Trace of a Matrix}

\begin{definitionblock}{Definition}
For a square matrix $A = (a_{ij}) \in \mathcal{M}_n ( \mathbb{R} )$, the \textbf{trace} of $A$, denoted $\operatorname{tr}(A)$, is defined by
\vspace{-5mm}
\[
\operatorname{tr}(A) = \sum_{i=1}^{n} a_{ii}.
\vspace{-2mm}
\]
That is, it is the sum of the diagonal elements of $A$.
\end{definitionblock}

\pause
\vspace{-5mm}
\begin{alertblock}{Key Properties}
For all $A,B \in \mathcal{M}_n ( \mathbb{R})$ and $\lambda \in \mathbb{R}$:
\vspace{-4mm}
\begin{itemize}
  \item $\operatorname{tr}(A+B) = \operatorname{tr}(A) + \operatorname{tr}(B)$
  \item $\operatorname{tr}(\lambda A) = \lambda\,\operatorname{tr}(A)$
  \item $\operatorname{tr}(A^\top) = \operatorname{tr}(A)$
  \item $\operatorname{tr}(AB) = \operatorname{tr}(BA)$ 
  \item $\operatorname{tr}(I_n) = n$
\end{itemize}
\end{alertblock}


\end{frame}

\begin{frame}{Examples of Trace Computation}

\begin{exampleblock}{Example 1 — Simple Numerical Matrix}
\[
A =
\begin{pmatrix}
2 & 5 & -1\\
0 & 3 & 4\\
7 & 1 & 1
\end{pmatrix}
\quad\Longrightarrow\quad
\operatorname{tr}(A) = 2 + 3 + 1 = 6.
\]
\end{exampleblock}

\pause

\begin{exampleblock}{Example 2 — Product of Matrices}
Let
$
A =
\begin{pmatrix}
1 & 2\\
0 & 3
\end{pmatrix},
\quad
B =
\begin{pmatrix}
4 & 1\\
2 & 0
\end{pmatrix}.
$
Then
\[
AB =
\begin{pmatrix}
1\cdot4 + 2\cdot2 & 1\cdot1 + 2\cdot0\\
0\cdot4 + 3\cdot2 & 0\cdot1 + 3\cdot0
\end{pmatrix}
=
\begin{pmatrix}
8 & 1\\
6 & 0
\end{pmatrix},
\quad
\operatorname{tr}(AB)=8+0=8.
\]
\[
BA =
\begin{pmatrix}
4\cdot1 + 1\cdot0 & 4\cdot2 + 1\cdot3\\
2\cdot1 + 0\cdot0 & 2\cdot2 + 0\cdot3
\end{pmatrix}
=
\begin{pmatrix}
4 & 11\\
2 & 4
\end{pmatrix},
\quad
\operatorname{tr}(BA)=4+4=8.
\]
\end{exampleblock}
\end{frame}

\begin{frame}{Examples of Trace Computation}

\begin{exampleblock}{Example 3 — Orthogonal Matrices}
If $Q \in O(n)$ (so $Q^\top Q = I_n$), then
\[
\operatorname{tr}(Q^\top Q) = \operatorname{tr}(I_n) = n.
\]
\end{exampleblock}

\end{frame}

\subsection{Determinant}
\begin{frame}{Determinant of a Square Matrix}

\begin{definitionblock}{Definition (Recursive Definition)}
For $A=(a_{ij})\in \mathcal{M}_{n} (\mathbb{R} )$, the \textbf{determinant} of $A$, denoted $\det(A)$, is defined recursively by:
\vspace{-4mm}
\[
\det(A) =
\begin{cases}
a_{11}, & n=1,\\[4pt]
\displaystyle\sum_{j=1}^{n} (-1)^{1+j} a_{1j} \det(A_{1j}), & n>1,
\end{cases}
\vspace{-3mm}
\]
where $A_{1j}$ is the submatrix obtained by deleting row 1 and column $j$.
\end{definitionblock}

\end{frame}

\begin{frame}{Explicit Determinant Formulas}

\begin{exampleblock}{For $n=2$}
If 
\[
A =
\begin{pmatrix}
a & b\\
c & d
\end{pmatrix},
\quad
\text{then} \quad
\det(A) = ad - bc.
\]
\end{exampleblock}

\pause

\begin{exampleblock}{For $n=3$}
If 
\[
A =
\begin{pmatrix}
a & b & c\\
d & e & f\\
g & h & i
\end{pmatrix},
\]
then by cofactor expansion:
\[
\det(A) = 
a(ei - fh) - b(di - fg) + c(dh - eg).
\]
\end{exampleblock}

\end{frame}

\begin{frame}{Fundamental Properties of the Determinant}

\begin{alertblock}{Key Properties}
For $A,B\in\mathcal{M}_{n} (\mathbb{R} )$ and $\lambda\in\mathbb{R}$:
\begin{itemize}
  \item $\det(I_n)=1$, \quad $\det(A^\top)=\det(A)$
  \item $\det(AB)=\det(A)\det(B)$
  \item $\det(\lambda A)=\lambda^n \det(A)$
  \item If $A$ has a zero row or two equal rows $\Rightarrow \det(A)=0$
  \item Swapping two rows multiplies $\det(A)$ by $-1$
  \item 
$
A\ \text{is invertible} 
\quad\Longleftrightarrow\quad
\det(A)\neq0
$ and if $A$ is invertible, $\det(A^{-1})= 1/\det(A)$.
\end{itemize}
\end{alertblock}
\end{frame}

\begin{frame}{Expansion by Row or Column}

\begin{definitionblock}{Cofactor Expansion Formula}
For any $A=(a_{ij})\in\mathcal{M}_{n} (\mathbb{R} )$ and any fixed row $i$ or column $j$:
\[
\det(A) = \sum_{k=1}^{n} (-1)^{i+k} a_{ik} \det(A_{ik})
= \sum_{k=1}^{n} (-1)^{k+j} a_{kj} \det(A_{kj}).
\]
\end{definitionblock}

\pause

\begin{exampleblock}{Example (expansion along the 2nd row)}
\[
A =
\begin{pmatrix}
1 & 2 & 3\\
4 & 0 & 6\\
7 & 8 & 9
\end{pmatrix}
\Rightarrow
\det(A)
= \sum_{j=1}^3 (-1)^{2+j} a_{2j}\, \det(A_{2j})
\]
\[
\begin{aligned}
\det(A)
&= (-1)^{2+1}\, \color{blue}{4}\,
   \det\!\begin{pmatrix} 2 & 3\\ 8 & 9 \end{pmatrix}
 \;+\; (-1)^{2+2}\, \color{blue}{0}\, (\cdot)
 \;+\; (-1)^{2+3}\, \color{blue}{6}\,
   \det\!\begin{pmatrix} 1 & 2\\ 7 & 8 \end{pmatrix} \\[4pt]
&= -4 \,\Big( 2\cdot 9 - 3\cdot 8 \Big)
   \;+\; 0
   \;-\; 6 \,\Big( 1\cdot 8 - 2\cdot 7 \Big) \\[4pt]
&= -4 \,\Big( 18 - 24 \Big)
   \;-\; 6 \,\Big( 8 - 14 \Big) \\[4pt]
&= -4 \,(-6) \;-\; 6 \,(-6) \\[4pt]
&= 24 \;+\; 36 \\[2pt]
&= \boxed{60}.
\end{aligned}
\]
\end{exampleblock}

\end{frame}

\begin{frame}{Comatrix and the Adjugate}

\begin{definitionblock}{Definition (Cofactor and Comatrix)}
For $A=(a_{ij})\in\mathcal{M}_{n} (\mathbb{R} )$, define the \textbf{cofactor}:
\vspace{-3.5mm}
\[
C_{ij} = (-1)^{i+j}\det(A_{ij}),
\vspace{-3mm}
\]
and the \textbf{comatrix} (or cofactor matrix)
\vspace{-3.5mm}
\[
\operatorname{Com}(A) = (C_{ij})_{i,j}.
\vspace{-3.5mm}
\]
The \textbf{adjugate} (or adjoint) matrix is 
\vspace{-3.5mm}
\[
\operatorname{adj}(A) = \operatorname{Com}(A)^\top.
\vspace{-3.5mm}
\]
\end{definitionblock}
\vspace{-5mm}
\pause

\begin{block}{Inverse Formula via the Adjugate}
If $\det(A)\neq 0$, then
\[
A^{-1} = \frac{1}{\det(A)}\,\operatorname{adj}(A).
\]
\textbf{Check:} $A\operatorname{adj}(A) = \operatorname{adj}(A)A = \det(A)I_n$.
\end{block}

\end{frame}

\begin{frame}{Summary: Determinant and Invertibility}

\begin{block}{What the Determinant Tells Us}
\begin{itemize}
  \item $\det(A) = 0$  $\Rightarrow$  $A$ is singular (not invertible), columns are linearly dependent.
  \item $\det(A) \neq 0$  $\Rightarrow$  $A$ is invertible, and the linear transformation preserves volume up to a scaling by $|\det(A)|$.
\end{itemize}
\end{block}

\pause

\begin{exampleblock}{Geometric Intuition (in $\mathbb{R}^2$)}
If $A$ maps the unit square to a parallelogram of area $|\det(A)|$:
\[
A =
\begin{pmatrix}
2 & 0\\
1 & 1
\end{pmatrix}
\Rightarrow
\det(A)=2,
\]
the area is doubled.
\end{exampleblock}

\pause

\textbf{Summary Formula:}
\[
A^{-1} = \frac{1}{\det(A)}\operatorname{adj}(A)
\quad\text{iff}\quad \det(A)\neq 0.
\]
\end{frame}

\begin{frame}{Application : inverse of a $2\time 2$ Matrix}

\begin{definitionblock}{Formula}
For
$
A =
\begin{pmatrix}
a & b\\
c & d
\end{pmatrix}
\quad\text{with}\quad \det(A)=ad-bc \neq 0,
$
the inverse is
\vspace{-3mm}
\[
A^{-1} = \frac{1}{\det(A)}
\begin{pmatrix}
d & -b\\
-c & a
\end{pmatrix}.
\vspace{-6mm}
\]
\end{definitionblock}

\pause
\vspace{-3mm}
\begin{exampleblock}{Example}
Let $
A =
\begin{pmatrix}
2 & 1\\
3 & 4
\end{pmatrix},
\quad
\det(A) = 2\times4 - 1\times3 = 5.
$
Then
$
A^{-1} = \frac{1}{5}
\begin{pmatrix}
4 & -1\\
-3 & 2
\end{pmatrix}
=
\begin{pmatrix}
0.8 & -0.2\\
-0.6 & 0.4
\end{pmatrix}.
$

\pause
\textbf{Verification:}
$
AA^{-1} =
\begin{pmatrix}
1 & 0\\
0 & 1
\end{pmatrix} = I_2.
$
\end{exampleblock}
    
\end{frame}

\section{Diagonalization}
  \begin{frame}{Outline}
    \tableofcontents[currentsection, hideothersubsections]
  \end{frame}
  
\subsection{Motivation}

\begin{frame}{What Does It Mean to Be Diagonalizable?}

\begin{definitionblock}{Definition}
\vspace{1mm}
A matrix $A \in \mathcal{M}_{n} (\mathbb{R} )$ is said to be \textbf{diagonalizable} if there exists an invertible matrix $P$ and a diagonal matrix $D$ such that
\vspace{-4mm}
\[
A = P D P^{-1}.
\vspace{-6mm}
\]
\end{definitionblock}
\vspace{-4mm}

\pause

\begin{block}{Interpretation}
\vspace{-2mm}
\begin{itemize}
  \item The columns of $P$ are called eigenvectors of $A$.
  \item The diagonal entries of $D$ are the corresponding eigenvalues.
  \item $A$ acts as a \textbf{scaling operator} on eigenvectors: $
  A v_i = \lambda_i v_i$.
\end{itemize}
\end{block}

\pause

\begin{exampleblock}{Example (Simple Case)}
If
\[
A =
\begin{pmatrix}
2 & 0\\
0 & 3
\end{pmatrix},
\quad \text{then} \quad
P = I_2,\ D = A,
\]
so $A$ is already diagonal — it stretches the first axis by $2$ and the second by $3$.
\end{exampleblock}

\end{frame}

\begin{frame}{Motivation: Why Diagonalize a Matrix? Principal Component Analysis (PCA)?}

High-dimensional datasets can bedifficult to visualize, computationally expensive to process, noisy and redundant (correlated variables).

  \pause
\begin{itemize}
  \item PCA finds a new orthogonal coordinate system such that:
  \begin{itemize}
    \item the first axis captures the maximum variance,
    \item the following axes capture decreasing variance,
    \item variables become uncorrelated.
  \end{itemize}
  \pause

  \item Mathematically, PCA diagonalizes the covariance matrix:
$
  C = P \Lambda P^{-1},
$
  where columns of $P$ are eigenvectors (principal directions)
  and $\Lambda$ contains eigenvalues (explained variances).
  \pause

  \item Dimensionality reduction:
  \vspace{-2mm}
  \[
  \text{Keep the first $k$ largest eigenvalues} \quad \Rightarrow \quad
  \text{project data onto $k$ principal components.}
  \]
    \vspace{-9mm}
\end{itemize}

\pause
\begin{exampleblock}{Example}
A dataset with strongly correlated variables in $\mathbb{R}^2$
can be represented accurately by projecting onto one principal axis (1D),
while preserving most of the variance.
\end{exampleblock}

\end{frame}


\begin{frame}{Motivation: Why Diagonalize a Matrix?}

\begin{itemize}
  \item Computing powers of a matrix $A^n$ is often required in:
  \begin{itemize}
    \item Solving discrete-time linear systems: $x_{t+1}=A x_t$.
    \item Studying Markov chains or stability of dynamical systems.
    \item Evaluating matrix exponentials $e^{At}$ in continuous-time models.
  \end{itemize}
  \pause
  \item If $A$ is diagonalizable, computations simplify:
  \[
  A = P D P^{-1} \quad \Rightarrow \quad A^n = P D^n P^{-1},
  \]
  where $D$ is diagonal.
  \item This avoids costly repeated multiplications.
\end{itemize}

\pause
\begin{exampleblock}{Example}
If $A = \begin{pmatrix}2 & 0\\0 & 3\end{pmatrix}$,  
then $A^n = \begin{pmatrix}2^n & 0\\0 & 3^n\end{pmatrix}$.
\end{exampleblock}

\end{frame}

\begin{frame}{Example: A Two-Dimensional Linear Recurrence System}

\begin{block}{System Definition}
Consider the coupled recurrences:
$
\begin{cases}
x_{n+1} = 2x_n + y_n,\\
y_{n+1} = x_n + y_n.
\end{cases}
$
\end{block}
\vspace{-5mm}
\pause

\begin{block}{Matrix Form}
We can write this system as
\[
\begin{pmatrix}
x_{n+1}\\[3pt]
y_{n+1}
\end{pmatrix}
=
\underbrace{
\begin{pmatrix}
2 & 1\\
1 & 1
\end{pmatrix}}_{A}
\begin{pmatrix}
x_n\\[3pt]
y_n
\end{pmatrix},
\quad \Rightarrow \quad
\begin{pmatrix}
x_n\\
y_n
\end{pmatrix}
= A^n
\begin{pmatrix}
x_0\\
y_0
\end{pmatrix}.
\]
\end{block}

\pause
\vspace{-6mm}

\begin{exampleblock}{Diagonalization Insight}
If $A = P D P^{-1}$ with eigenvalues $\lambda_1, \lambda_2$, then
\[
A^n = P D^n P^{-1},
\quad \text{so} \quad
D^n =
\begin{pmatrix}
\lambda_1^n & 0\\
0 & \lambda_2^n
\end{pmatrix}.
\]
Hence each component evolves geometrically at rate $\lambda_i^n$.
\end{exampleblock}

\end{frame}

\begin{frame}{Example: A Second-Order Recurrence as a First-Order System}

\begin{block}{Scalar Recurrence}
Consider the sequence $(u_n)$ defined by
\vspace{-5mm}
\[
u_{n+2} = 3u_{n+1} - 2u_n, \quad u_0 = 1, \quad u_1 = 2.
\vspace{-6mm}
\]
\end{block}

\pause

\begin{block}{Rewriting as a Matrix Recurrence}
Define the vector
$
X_n =
\begin{pmatrix}
u_{n+1}\\
u_n
\end{pmatrix}.
$
Then
$
X_{n+1} =
\begin{pmatrix}
u_{n+2}\\
u_{n+1}
\end{pmatrix}
=
\underbrace{
\begin{pmatrix}
3 & -2\\
1 & 0
\end{pmatrix}}_{A}
\begin{pmatrix}
u_{n+1}\\
u_n
\end{pmatrix}
= A X_n.
$
Hence
$
X_n = A^n X_0, \quad \text{with } X_0 = 
\begin{pmatrix}u_1\\u_0\end{pmatrix}
= 
\begin{pmatrix}2\\1\end{pmatrix}.
$
\end{block}

\pause

\begin{exampleblock}{Conclusion}
The scalar recurrence $u_{n+2}=3u_{n+1}-2u_n$ is equivalent to the matrix recurrence $X_{n+1}=A X_n$.  
Once diagonalized, $A^n$ gives a closed-form expression for $u_n$.
\end{exampleblock}

\end{frame}


\subsection{Characteristic Polynomial}

\begin{frame}{Characteristic Polynomial}

\begin{definitionblock}{Definition}
For $A \in \mathcal{M}_n (\mathbb{R})$, the \textbf{characteristic polynomial} of $A$ is
\[
p_A(\lambda) = \det(A - \lambda I_n).
\]
The \textbf{eigenvalues} of $A$ are the roots of $p_A(\lambda)=0$.
\end{definitionblock}

\pause
\vspace{-4mm}

\begin{exampleblock}{Example: $2\times2$ Matrix}
\[
A = 
\begin{pmatrix}
2 & 1\\
1 & 2
\end{pmatrix}.
\]
Then
\[
p_A(\lambda)
= \det
\begin{pmatrix}
2-\lambda & 1\\
1 & 2-\lambda
\end{pmatrix}
= (2-\lambda)^2 - 1 = \lambda^2 - 4\lambda + 3.
\]
Hence, $\lambda_1=1$ and $\lambda_2=3$.
\end{exampleblock}
\end{frame}

\begin{frame}{Eigenvectors and Geometric Multiplicity}

\begin{definitionblock}{Definition}
For $A \in \mathcal{M}_n (\mathbb{R})$, the \textbf{characteristic polynomial} of $A$ is
\[
p_A(\lambda) = \det(A - \lambda I_n).
\]
The \textbf{eigenvalues} of $A$ are the roots of $p_A(\lambda)=0$.
\end{definitionblock}

\begin{definitionblock}{Eigenvectors}
A nonzero vector $x \in \mathbb{R}^n$ is an \textbf{eigenvector} of $A$ associated to eigenvalue $\lambda$ if
\[
A x = \lambda x.
\]
\end{definitionblock}
\end{frame}

\subsection{Algorithmic Methods to diagonalize}

\begin{frame}{Algorithm to Diagonalize a Matrix}

\begin{enumerate}
  \item Compute the characteristic polynomial $p_A(\lambda) = \det(A - \lambda I)$.
  \item Find all eigenvalues $\lambda_1,\ldots,\lambda_k$ (solve $p_A(\lambda)=0$).
  \item For each $\lambda_i$, find eigenvectors by solving
  \[
  (A - \lambda_i I)x = 0.
  \]
  \item Form the matrix $P$ whose columns are eigenvectors.
  \item Then $D = P^{-1} A P$ is diagonal, with eigenvalues on the diagonal.
\end{enumerate}

\pause
\begin{block}{Check}
$A$ is diagonalizable $\Leftrightarrow$ the matrix $P$ of eigenvectors is invertible.
\end{block}

\end{frame}

\begin{frame}{Example: Diagonalization of a $2\times2$ Matrix}

\begin{itemize}
  \item Compute the characteristic polynomial:
  \[
  p_A(\lambda) = \det(A - \lambda I) 
  = (4-\lambda)(3-\lambda) - 2 
  = \lambda^2 - 7\lambda + 10.
  \]
  \item Roots (eigenvalues):
  \[
  \textcolor{blue}{\lambda_1 = 5}, 
  \qquad 
  \textcolor{red}{\lambda_2 = 2.}
  \]

  \pause

  \item Eigenvectors:
  \[
  (A - \textcolor{blue}{5}I)x = 0 
  \Rightarrow 
  \begin{pmatrix}
  -1 & 1\\
  2 & -2
  \end{pmatrix}
  x = 0
  \Rightarrow 
  \textcolor{blue}{v_1 = 
  \begin{pmatrix}1\\1\end{pmatrix}}.
  \]
  \[
  (A - \textcolor{red}{2}I)x = 0 
  \Rightarrow 
  \begin{pmatrix}
  2 & 1\\
  2 & 1
  \end{pmatrix}
  x = 0
  \Rightarrow 
  \textcolor{red}{v_2 = 
  \begin{pmatrix}1\\-2\end{pmatrix}}.
  \]
\end{itemize}

\pause
\[
P = 
\begin{pmatrix}
\textcolor{blue}{1} & \textcolor{red}{1}\\
\textcolor{blue}{1} & \textcolor{red}{-2}
\end{pmatrix},
\quad
D = 
\begin{pmatrix}
\textcolor{blue}{5} & 0\\
0 & \textcolor{red}{2}
\end{pmatrix},
\quad
A = P D P^{-1}.
\]
\end{frame}

\begin{frame}{Verifying $A = P D P^{-1}$ (Explicit Computation)}

We claimed that 
\[
A =
\begin{pmatrix}
4 & 1\\
2 & 3
\end{pmatrix}
=
P D P^{-1}
\quad\text{with}\quad
P =
\begin{pmatrix}
1 & 1\\
1 & -2
\end{pmatrix},
\quad
D =
\begin{pmatrix}
5 & 0\\
0 & 2
\end{pmatrix}.
\]

\pause
\begin{block}{Step 1: Compute $P^{-1}$}
\vspace{2mm}

For a $2\times 2$ matrix 
$\begin{pmatrix}a & b\\ c & d\end{pmatrix}$,
$
\begin{pmatrix}a & b\\ c & d\end{pmatrix}^{-1}
= \frac{1}{ad-bc}
\begin{pmatrix}
d & -b\\
-c & a
\end{pmatrix}.
$


\vspace{1mm}
\noindent 
Here $a=1$, $b=1$, $c=1$, $d=-2$, so
$
\det(P)= (1)(-2) - (1)(1) = -3,
$
\[
P^{-1}
= \frac{1}{-3}
\begin{pmatrix}
-2 & -1\\
-1 & 1
\end{pmatrix}
=
\begin{pmatrix}
\frac{2}{3} & \frac{1}{3}\\[4pt]
\frac{1}{3} & -\frac{1}{3}
\end{pmatrix}.
\]
\end{block}
\end{frame}

\begin{frame}{Verifying $A = P D P^{-1}$ (Explicit Computation)}
\vspace{-5mm}
\[
A =
\begin{pmatrix}
4 & 1\\
2 & 3
\end{pmatrix}
=
P D P^{-1}
\,\, \text{with} \,\,
P =
\begin{pmatrix}
1 & 1\\
1 & -2
\end{pmatrix},
\,\,
D =
\begin{pmatrix}
5 & 0\\
0 & 2
\end{pmatrix},
\,\,
P^{-1}
=\frac{1}{3}
\begin{pmatrix}
2 & 1\\
1 & -1
\end{pmatrix}.
\vspace{-5mm}
\]
\begin{block}{Step 2: Compute $P D$}
\vspace{-4mm}
\[
P D
=
\begin{pmatrix}
1 & 1\\
1 & -2
\end{pmatrix}
\begin{pmatrix}
5 & 0\\
0 & 2
\end{pmatrix}
=
\begin{pmatrix}
1\cdot5 + 1\cdot0 & 1\cdot0 + 1\cdot2\\
1\cdot5 + (-2)\cdot0 & 1\cdot0 + (-2)\cdot2
\end{pmatrix}
=
\begin{pmatrix}
5 & 2\\
5 & -4
\end{pmatrix}.
\vspace{-6mm}
\]
\end{block}

\pause
\begin{block}{Step 3: Compute $P D P^{-1}$}
Now multiply:
\begin{align*}
P D P^{-1}
= \frac{1}{3}
\begin{pmatrix}
5 & 2\\
5 & -4
\end{pmatrix}
\begin{pmatrix}
2& 1\\
1 & -1
\end{pmatrix}
&= \frac{1}{3}
\begin{pmatrix}
5\cdot2 + 2\cdot1 & 5\cdot1 + 2\cdot(-1)\\[4pt]
5\cdot2 + (-4)\cdot1 & 5\cdot1 + (-4)\cdot(-1)
\end{pmatrix}
\\
&= \frac{1}{3}
\begin{pmatrix}
12 & 3\\
6 & 9
\end{pmatrix}
=
\begin{pmatrix}
4 & 1\\
2 & 3
\end{pmatrix}
=
A
\end{align*}
\end{block}

\end{frame}

\begin{frame}{Carefull ! Not all matrix are diagonalizable}

Let
$
A =
\begin{pmatrix}
0 & 1\\
0 & 0
\end{pmatrix}.
$
Note that
$
A^2 =
\begin{pmatrix}
0 & 1\\
0 & 0
\end{pmatrix}
\begin{pmatrix}
0 & 1\\
0 & 0
\end{pmatrix}
=
\begin{pmatrix}
0 & 0\\
0 & 0
\end{pmatrix}
= 0.
$
Such a matrix is called \textbf{nilpotent} (here: $A^2 = 0$ but $A \neq 0$).
\pause
We prove that $A$ is \textbf{not diagonalizable} by contradiction.
\pause
\medskip
Assume there exists an invertible $P$ and a diagonal matrix $D$ such that
$
A = P D P^{-1}.
$
Then
$
A^2 = (P D P^{-1})(P D P^{-1}) = P D^2 P^{-1}.
$

But we already know $A^2 = 0$, so
$
0 = A^2 = P D^2 P^{-1}
\quad \Longrightarrow \quad
D^2 = 0.
$

\pause

\begin{block}{Contradiction}
If $D$ is diagonal and $D^2 = 0$, then every diagonal entry of $D$ must be $0$.
So
\[
D = 0
\quad \Longrightarrow \quad
A = P D P^{-1} = P (0) P^{-1} = 0.
\]

But $A \neq 0$:
\[
A =
\begin{pmatrix}
0 & 1\\
0 & 0
\end{pmatrix}
\neq 0.
\]

\[
\text{Contradiction.}
\]

Therefore $A$ cannot be diagonalizable.
\end{block}

\end{frame}

\subsection{A special case : Symmetric Matrices}

\begin{frame}{Spectral Theorem for Symmetric Matrices}

\begin{theorem}[Spectral Theorem]
If $A \in \mathcal{M}_{n} (\mathbb{R} )$ is \textbf{symmetric} ($A^\top = A$), then:
\begin{itemize}
  \item All eigenvalues of $A$ are real.
  \item There exists an orthogonal matrix $P$ (\textit{i.e.} $P P^\top = I_n$) and $D$ diagonal, such that
  \[
  A = P D P^\top,
  \]
\end{itemize}
\end{theorem}

\pause
\begin{block}{Interpretation}
\begin{itemize}
  \item Important in Principal Component Analysis (PCA), covariance matrices, and risk models.
\end{itemize}
\end{block}

\end{frame}

\begin{frame}{Example: Diagonalization of a Symmetric Matrix}

\[
A =
\begin{pmatrix}
2 & 1\\
1 & 2
\end{pmatrix}.
\]
\pause
\begin{itemize}
  \item $p_A(\lambda) = \det(A - \lambda I) = (2-\lambda)^2 - 1 = \lambda^2 - 4\lambda + 3.$
  \item Eigenvalues: $\lambda_1=1$, $\lambda_2=3.$
  \item Eigenvectors:
  \[
  v_1 = \frac{1}{\sqrt{2}}
  \begin{pmatrix}1\\-1\end{pmatrix},
  \quad
  v_2 = \frac{1}{\sqrt{2}}
  \begin{pmatrix}1\\1\end{pmatrix}.
  \]
  \item Orthonormal basis: $P = [v_1 \ v_2]$.
\end{itemize}

\pause
\[
P^\top A P = 
\begin{pmatrix}
1 & 0\\
0 & 3
\end{pmatrix}.
\]
\textbf{Conclusion:} $A$ is symmetric $\Rightarrow$ diagonalizable by an orthogonal matrix.
\end{frame}

\section{Exponential Matrices}

\begin{frame}{Norms on Matrices}
\begin{definitionblock}{Definition}
A \textbf{matrix norm} $\|\cdot\|$ is a function $\mathcal{M}_{n} (\mathbb{R} )\to\mathbb{R}_+$ such that, for all $A,B\in\mathcal{M}_{n} (\mathbb{R} )$ and $\lambda\in\mathbb{R}$:
\begin{itemize}
    \item $\|A\| = 0 \iff A=0$,
    \item $\|\lambda A\| = |\lambda|\,\|A\|$,
    \item $\|A+B\| \le \|A\| + \|B\|$ (triangle inequality),
    \item $\|AB\| \le \|A\|\,\|B\|$ (sub-multiplicativity).
\end{itemize}
\end{definitionblock}
\pause
\begin{alertblock}{Reminder}
\begin{itemize}
    \item The first three points constitute the usual definition of a norm
    \item All norms on a finite-dimensional vector space are \textbf{equivalent}, which means that the norm you choose is not that important
\end{itemize}
\end{alertblock}
\end{frame}

\begin{frame}{Normal Convergence of a Series}
\begin{definitionblock}{Definition}
A series $\sum_{k\ge0} A_k$ of matrices is said to be \textbf{normally convergent} if  
$ \lim_{n \rightarrow \infty} \displaystyle \sum^n_{k=0} \|A_k\| < \infty$.
\end{definitionblock}
\pause
\begin{block}{Theorem}
If $\sum A_k$ converges normally, then it converges (absolutely) in $\mathcal{M}_{n} (\mathbb{R} )$, and the limit is independent of the chosen norm.
\end{block}
\end{frame}

\begin{frame}{Exponential of a Matrix}
\begin{definitionblock}{Definition}
For $A\in\mathbb{R}^{n\times n}$, the \textbf{matrix exponential} is defined by the convergent power series:
\[
e^{A} = \sum_{k=0}^{\infty} \frac{A^k}{k!}.
\]
\end{definitionblock}
\pause
\begin{exampleblock}{Example (Nilpotent Matrix)}
If $N=\begin{pmatrix}0&1\\0&0\end{pmatrix}$, then $N^2=0$.  
Hence:
\[
e^{N} = I + N = 
\begin{pmatrix}
1 & 1\\
0 & 1
\end{pmatrix}.
\]
\end{exampleblock}
\end{frame}

\begin{frame}{Key Properties of the Matrix Exponential}
\begin{itemize}
    \item $e^{0}=I$, \quad $(e^{A})^{-1} = e^{-A}$.
    \item $\det(e^{A}) = e^{\operatorname{tr}(A)}$.
    \item If $A$ and $B$ commute ($AB=BA$), then $e^{A+B}=e^{A}e^{B}$.
    \item If $A$ is diagonalizable ($A=PDP^{-1}$):
    \[
    e^{A} = P\,e^{D}\,P^{-1}, \quad e^{D}=\operatorname{diag}(e^{\lambda_1},\ldots,e^{\lambda_n}).
    \]
\end{itemize}
\pause
\begin{alertblock}{Spectral Link}
Eigenvalues of $e^{A}$ are $e^{\lambda_i}$, where $\lambda_i$ are the eigenvalues of $A$.
\end{alertblock}
\end{frame}

\begin{frame}{Be Careful!}
\begin{alertblock}{Common Misconceptions}
\begin{itemize}
    \item In general, $e^{A+B} \ne e^{A} e^{B}$ if $AB \ne BA$.
    \item $e^{A}$ is not obtained by exponentiating each entry of $A$.
\end{itemize}
\end{alertblock}
\pause
\begin{exampleblock}{Counterexample}
Let $A=\begin{pmatrix}0&1\\0&0\end{pmatrix}$ and $B=\begin{pmatrix}0&0\\1&0\end{pmatrix}$.  
Then $AB\ne BA$ and $e^{A+B}\ne e^{A}e^{B}$.
\end{exampleblock}
\end{frame}

\begin{frame}{Time Derivative in Finite Dimension}
\begin{definitionblock}{Derivative of a Vector Function}
Let $f:\mathbb{R}\to\mathbb{R}^n$. We say $f$ is differentiable at $t_0$ if
\[
\lim_{h\to0}\frac{\|f(t_0+h)-f(t_0)-h f'(t_0)\|}{|h|}=0.
\]
\end{definitionblock}
\pause
\begin{block}{Norm Reminder}
Any norm on $\mathbb{R}^n$ induces the same notion of derivative.  
We can thus use the Euclidean norm safely.
\end{block}
\end{frame}

\begin{frame}{Differentiating an Infinite Sum}
\begin{alertblock}{Theorem (Term-by-Term Differentiation)}
Let $(f_k)_{k\ge0}$ be a sequence of $\mathcal{C}^1$ functions on an interval $I\subset\mathbb{R}$.  
If
\[
\sum_{k\ge0} f_k(t)
\quad\text{and}\quad
\sum_{k\ge0} f_k'(t)
\]
both converge \textbf{normally} (that is, $\sum \|f_k(t)\|$ and $\sum \|f_k'(t)\|$ converge uniformly on $I$),
then:
\[
\frac{d}{dt}\!\left(\sum_{k=0}^{\infty} f_k(t)\right)
= \sum_{k=0}^{\infty} f_k'(t),
\quad \text{for all } t\in I,
\]
and the convergence of both series is uniform on $I$.
\end{alertblock}

\pause
\begin{block}{Idea of the Proof}
Normal convergence allows us to exchange the limit defining the derivative
\[
\lim_{h\to0} \frac{S(t+h)-S(t)}{h}
\]
with the infinite sum $\sum f_k(t)$, because the tail terms can be uniformly controlled.
\end{block}

\pause
\begin{exampleblock}{Application: the exponential series}
For a fixed matrix $A\in\mathbb{R}^{n\times n}$,
\[
e^{tA} = \sum_{k=0}^{\infty} \frac{t^k A^k}{k!}
\]
converges normally on $\mathbb{R}$ since
\[
\sum_{k=0}^{\infty} \frac{|t|^k \|A\|^k}{k!} = e^{|t|\|A\|} < \infty.
\]
Thus we can differentiate term by term:
\[
\frac{d}{dt} e^{tA}
= \sum_{k=1}^{\infty} \frac{k\,t^{k-1}A^k}{k!}
= \sum_{k=1}^{\infty} \frac{t^{k-1}A^k}{(k-1)!}
= A \sum_{k=0}^{\infty} \frac{t^kA^k}{k!}
= A e^{tA}.
\]
\end{exampleblock}
\end{frame}
%======================================================
\begin{frame}{Time Derivative of a Matrix-Valued Function}
\begin{definitionblock}{Definition}
Let $A:\mathbb{R}\to\mathbb{R}^{n\times n}$ be a function.  
The \textbf{derivative} of $A$ at $t$ is the matrix
\[
A'(t) = \left( \frac{d}{dt} A_{ij}(t) \right)_{1\le i,j\le n}.
\]
\end{definitionblock}
\pause
\begin{exampleblock}{Example}
Let
\[
A(t) =
\begin{pmatrix}
e^t & t^2\\
0 & \sin t
\end{pmatrix}.
\]
Then
\[
A'(t) =
\begin{pmatrix}
e^t & 2t\\
0 & \cos t
\end{pmatrix}.
\]
\end{exampleblock}
\pause
\begin{block}{Observation}
If $A(t)=e^{tM}$ for a constant matrix $M$,  
then by term-by-term differentiation:
\[
A'(t) = M e^{tM}.
\]
\end{block}
\end{frame}

%======================================================
\begin{frame}{Linear ODE with Constant Coefficients}
\begin{definitionblock}{Theorem (Existence and Uniqueness)}
Let $A\in\mathbb{R}^{n\times n}$ and $b\in\mathbb{R}^n$.  
For the system
\[
\dot{x}(t) = A x(t) + b, \qquad x(0)=x_0,
\]
there exists a unique continuously differentiable solution given by:
\[
x(t) = e^{tA}x_0 + \int_0^t e^{(t-s)A}b\,ds.
\]
\end{definitionblock}
\pause
\begin{exampleblock}{Homogeneous Case ($b=0$)}
If $b=0$, then
\[
x(t) = e^{tA}x_0.
\]
Each component of $x(t)$ is a linear combination of the exponentials of $A$'s eigenvalues.
\end{exampleblock}
\end{frame}

%======================================================
\begin{frame}{Second-Order ODE as a First-Order System}
\begin{definitionblock}{Reduction}
Consider a second-order vector equation:
\[
x''(t) = A x(t), \quad x(t)\in\mathbb{R}^n.
\]
Define
\[
Y(t) =
\begin{pmatrix}
x(t)\\
x'(t)
\end{pmatrix}.
\]
Then
\[
Y'(t) =
\begin{pmatrix}
0 & I_n\\
A & 0
\end{pmatrix}
Y(t).
\]
\end{definitionblock}
\pause
\begin{block}{Interpretation}
The dynamics of an order-2 system can be represented as a linear system of dimension $2n$ with a constant block matrix.  
This allows the use of the matrix exponential:
\[
Y(t) = e^{tB} Y(0), \quad B=\begin{pmatrix}0&I\\A&0\end{pmatrix}.
\]
\end{block}
\end{frame}

%======================================================
\begin{frame}{Example: Solving a Linear System via the Matrix Exponential}
\begin{exampleblock}{Example}
Consider
\[
\dot{x}(t) =
\begin{pmatrix}
-2 & 1\\
0 & -1
\end{pmatrix}
x(t),
\qquad
x(0)=
\begin{pmatrix}
1\\
1
\end{pmatrix}.
\]
\begin{enumerate}
\item The matrix is upper triangular; its eigenvalues are $\lambda_1=-2$ and $\lambda_2=-1$.
\item It is already in a nearly diagonal form, so $e^{tA}$ is easily computed:
\[
e^{tA} =
\begin{pmatrix}
e^{-2t} & e^{-t}-e^{-2t}\\
0 & e^{-t}
\end{pmatrix}.
\]
\item The solution is then
\[
x(t) = e^{tA} x_0 =
\begin{pmatrix}
e^{-2t} + e^{-t} - e^{-2t}\\
e^{-t}
\end{pmatrix}
=
\begin{pmatrix}
e^{-t}\\ e^{-t}
\end{pmatrix}.
\]
\end{enumerate}
\end{exampleblock}
\pause
\begin{alertblock}{Long-Time Behavior}
Since $\operatorname{Re}(\lambda_i)<0$, all components decay exponentially:
\[
x(t)\to0 \quad \text{as } t\to\infty.
\]
\end{alertblock}
\end{frame}


\section{Numerical Linear Algebra in Python: NumPy and PyTorch}

%----------------------------------------------------------
\subsection{Why NumPy and PyTorch?}

\begin{frame}{Motivation}

\begin{itemize}
    \item Modern linear algebra is done on computers.
    \item Python provides two powerful libraries:
    \begin{itemize}
        \item \textbf{NumPy}: scientific computing and matrices.
        \item \textbf{PyTorch}: machine learning and automatic differentiation.
    \end{itemize}

    \pause

    \item Goals of this section:
    \begin{itemize}
        \item manipulate vectors and matrices,
        \item perform algebraic operations,
        \item use decomposition and optimization tools.
    \end{itemize}
\end{itemize}

\end{frame}

%----------------------------------------------------------
\subsection{NumPy Basics}

\begin{frame}{NumPy: Creating Arrays}

\begin{block}{Creation commands}
\begin{itemize}
    \item \texttt{np.array([[1,2],[3,4]])}
    \item \texttt{np.zeros((m,n))}
    \item \texttt{np.ones((m,n))}
    \item \texttt{np.eye(n)} (identity matrix)
\end{itemize}
\end{block}

\pause

\begin{alertblock}{Type}
NumPy arrays are called \textbf{ndarrays} (n-dimensional arrays).
\end{alertblock}

\end{frame}

%----------------------------------------------------------
\begin{frame}{NumPy: Basic Operations}

\begin{itemize}
    \item Scalar multiplication: \texttt{2*A}
    \item Addition: \texttt{A + B}
    \item Hadamard product (entrywise): \texttt{A * B}
    \item Matrix multiplication:
\[
\texttt{A @ B \quad or \quad np.matmul(A,B)}
\]
\end{itemize}

\pause

\begin{alertblock}{Reminder}
Matrix multiplication is not the same as entrywise multiplication.
\end{alertblock}

\end{frame}

%----------------------------------------------------------
\begin{frame}{NumPy: Linear Algebra Tools}

\begin{itemize}
    \item Transpose: $\text{A.T}$
    \item Inverse: $\text{np.linalg.inv(A)}$
    \item Determinant: $\text{np.linalg.det(A)}$
    \item Rank: $\text{np.linalg.matrix\_rank}(A)$
    \item Eigenvalues/eigenvectors:
$$
w, v = \text{np.linalg.eig}(A)
$$
\end{itemize}

\end{frame}

%----------------------------------------------------------
\subsection{PyTorch Basics}

\begin{frame}{PyTorch: What is it?}

\begin{block}{Main ideas}
\begin{itemize}
    \item library for numerical computation,
    \item designed for deep learning,
    \item optimized for GPU computation.
\end{itemize}
\end{block}

\pause

\begin{alertblock}{Key concept}
PyTorch uses \textbf{tensors}, a generalization of matrices.
\end{alertblock}

\end{frame}

%----------------------------------------------------------
\begin{frame}{PyTorch: Creating Tensors}

\begin{itemize}
    \item $\texttt{torch.tensor([[1.,2.],[3.,4.]])}$
    \item $\texttt{torch.zeros((m,n))}$
    \item $\texttt{torch.ones((m,n))}$
    \item $\texttt{torch.eye(n)}$
\end{itemize}

\pause

\begin{alertblock}{Note}
By default, PyTorch uses 32-bit floating point values.
\end{alertblock}

\end{frame}

%----------------------------------------------------------
\begin{frame}{PyTorch: Matrix Operations}

\begin{itemize}
    \item Scalar multiplication: \texttt{3*A}
    \item Addition: \texttt{A + B}
    \item Entrywise product: \texttt{A * B}
    \item Matrix multiplication:
$$
\texttt{A @ B \quad or \quad torch.matmul(A,B)}
$$
\end{itemize}

\end{frame}

%----------------------------------------------------------
\begin{frame}{PyTorch vs NumPy}

\begin{block}{Similarities}
\begin{itemize}
    \item Same syntax for most operations,
    \item Similar data structures.
\end{itemize}
\end{block}

\pause

\begin{alertblock}{Differences}
\begin{itemize}
    \item PyTorch supports \textbf{automatic differentiation},
    \item PyTorch runs on \textbf{GPU},
    \item NumPy is CPU only.
\end{itemize}
\end{alertblock}

\end{frame}

%----------------------------------------------------------
\subsection{Automatic Differentiation in PyTorch}

\begin{frame}{Automatic Differentiation}

\begin{itemize}
    \item For optimization, we need gradients.
    \item PyTorch computes them automatically.
\end{itemize}

\pause

\begin{block}{Example idea}
$$
\texttt{x = torch.tensor(2., requires\_grad=True)}
$$
PyTorch tracks operations applied to \texttt{x}
and can compute:
$$
\frac{d f}{dx}
$$
automatically.
\end{block}

\end{frame}

%----------------------------------------------------------
\begin{frame}{Gradient Computation Example}

\begin{itemize}
    \item Define a function:
$$
y = x^2 + 3x
$$
    \item Command to compute gradient:
$$
\texttt{y.backward()}
$$
    \item Retrieve gradient:
$$
\texttt{x.grad}
$$
\end{itemize}

\pause

\begin{alertblock}{Interpretation}
This is essential for:
\begin{itemize}
    \item regression,
    \item neural networks,
    \item any optimization problem.
\end{itemize}
\end{alertblock}

\end{frame}

%----------------------------------------------------------
\subsection{Summary}

\begin{frame}{Key Takeaways}

\begin{itemize}
    \item NumPy:
    \begin{itemize}
        \item scientific computing
        \item matrices and linear algebra
    \end{itemize}

    \pause

    \item PyTorch:
    \begin{itemize}
        \item tensors
        \item GPU support
        \item automatic differentiation
    \end{itemize}

    \pause

    \item Both are essential tools in modern applied mathematics.
\end{itemize}

\end{frame}

\end{document}